\section{Actual preordinary images in the ramified Artin coefficient}
\label{sec:preordinary-artin-image}

The applied primary source is Kings--Loeffler--Zerbes,
\emph{Rankin--Eisenstein classes and explicit reciprocity laws},
arXiv:1503.02888v5:
\url{https://discovery.ucl.ac.uk/id/eprint/1497960/1/Zerbes-S_1503.02888v5.pdf}.
All KLZ locators below refer to this version.

\begin{proposition}[AC-K33-image: the actual global classes and their coefficient]
\label{AC-K33-image}
Put \(p=269\), let \(e\geq1\), and let \(f\) be the fixed ordinary cuspidal
form of weight \(2e+1\) and nebentypus \(\eta_K\) occurring in the audited
Type~II representation.  We use the arithmetic-Frobenius normalization
\[
 \det(1-X\rho_f(\operatorname{Frob}_\ell))
  =1-a_\ell(f)X+\eta_K(\ell)\ell^{2e}X^2,\qquad
 V=\rho_f(-e),\qquad \det V=\eta_K.                 \tag{K33I.1}
\]
Let \(\rho_g=\operatorname{Ind}_{G_K}^{G_{\mathbf Q}}\chi\) be the cubic
Artin representation used in AC-K33.  Either cubic exponent
\(b\in\{1,2\}\) is allowed, with \(\rho_g=\operatorname{Ind}\chi^b\);
this \(b\) is distinct from the coefficient index \(j=e-1\).
A simultaneous quadratic
compensation \(f_0=f\otimes\tau,\ g_0=\theta_\chi\otimes\tau\) is allowed:
then use \(\rho_{f_0},\rho_{g_0}\) below, whose tensor product is the same
as \(\rho_f\otimes\rho_g\).  Fix the two-level curve and the simultaneous
coefficient quotient of AC-K33, after increasing the fixed tame level
and the fixed \(p\)-power level to accommodate \(f\).  All components
over \(\mathbf Q\) are retained.

For \(j=e-1\), the literal preordinary Rankin--Iwasawa classes, transported
by AC-K33-cyclotomic, have canonical images
\[
 {\cal C}_m\in
 H^1\bigl(\mathbf Q(\mu_m),
       \rho_f\otimes\rho_g^\vee(-j)\bigr)
       \otimes_E D^\vee,\qquad
 D=D_{\rm dR}\bigl(\rho_{g,p}(1)\bigr),             \tag{K33I.2}
\]
where \(D^\vee\) is a trivial global Galois multiplicity label,
for every auxiliary \(m\) prime to the fixed excluded primes.  Thus every
\(d\in D\) gives an actual class \(c_{m,d}\), linearly in \(d\).
These are obtained before either ordinary projection on the second
factor or inversion of its \(U_p'\).  A single lattice denominator
suffices for the fixed coefficient as \(m\) varies.

With an explicitly chosen induced quadratic self-twist of \(\rho_g\),
the coefficient in (K33I.2) is isomorphic to
\[
 U=(V\otimes\rho_g)^\vee(1).                       \tag{K33I.3}
\]
We may and do take \(3\mid N\).  Via Shapiro, these are also actual
classes over \(K(\mu_m)\) with the rank-two coefficient
\((V|_{G_K}\otimes\chi^b)^\vee(1)\), with the corresponding
local identification at the inert prime \(269\) for \(m=1\).

This proposition asserts construction and functoriality only.
It does not assert that any \(c_{m,d}\) is nonzero, that their local
singular images span a two-dimensional space, or that their cyclotomic
norms give the required Euler system over general \(K\)-ray class fields.
\end{proposition}

\begin{proof}
We specify the sheaves before taking cohomology.  In the notation of
KLZ, let
\[
 {\cal H}=(R^1\pi_*\mathcal O)^\vee,\qquad
 {\cal L}^{[j]}
 =\Lambda({\cal H}\langle t_N\rangle)
       \otimes_{\mathcal O}\operatorname{TSym}^j{\cal H}.
                                                               \tag{K33I.4}
\]
Definition~5.1.5 of KLZ places
\({}_c\mathrm{RI}^{[j]}_{M,N,a}\) in
\[
 H^3_{\rm et}\bigl(Y(M,N)^2,\,
       ({\cal L}^{[j]})^{\boxtimes2}(2-j)\bigr).    \tag{K33I.5}
\]
Here KLZ's square initially means the fibre product over its finite
cyclotomic base.  If we compute Kunneth on the full product over
\(\mathbf Q\), the same-component product is a Galois-stable clopen
union.  Extension by zero along this union is the canonical integral
map used below.  Thus retaining all \(\mathbf Q\)-components does not
replace the original class by a class on a chosen geometric component.
The arithmetic base is obtained by inverting the fixed bad primes.
For \(m\) prime to these primes and to \(c\), apply the exact two-level
transport of AC-K33-cyclotomic to the class at \((Mm,Nm)\); write the
result as \(z_m^{[j]}\).  It is defined over the full cyclotomic base
\(\mu_m^\circ\).  The transport uses \([m]\) on the measure factor and
\(\operatorname{TSym}^j(\phi)\) on the algebraic factor, precisely as
specified there.  In particular we introduce no extra \(m^j\) into
(K33I.5).

Write \(I_j=H^1_{\rm et}(Y(M,N)_{\overline{\mathbf Q}},
{\cal L}^{[j]}(1))\), interpreted as the actual compact inverse tower.
The two twists by \(1\) in the two copies of \(I_j\) account for the
\(2\) in (K33I.5); the remaining twist is exactly \(-j\).

Here are finite-coefficient and inverse-limit details of the
Hochschild--Serre map.  Set \(R_n=\mathcal O/p^n\mathcal O\), and first
fix an internal division-torsor level.  Each coefficient sheaf is
finite locally free over \(R_n\).  After an embedding into
\(\mathbf C\), a connected affine smooth curve has the homotopy type
of a finite graph.  The comparison theorem for finite coefficients
therefore represents its geometric cohomology by a two-term complex
\[
 C_i^0\xrightarrow{d_i}C_i^1
\]
of finite free \(R_n\)-modules.  This proves perfect amplitude
\([0,1]\), including for the finite free pushforwards of the division
torsors.  Derived Kunneth is represented by \(C_1^\bullet\otimes
C_2^\bullet\).  Its degree-two cohomology is, explicitly,
\[
 \frac{C_1^1\otimes C_2^1}
      {d_1C_1^0\otimes C_2^1+C_1^1\otimes d_2C_2^0}
   \;=\;H^1(C_1^\bullet)\otimes_{R_n}H^1(C_2^\bullet).
                                                               \tag{K33I.6}
\]
Thus no flatness assertion about the cohomology groups is needed.
There is no geometric cohomology above degree \(2\) on the product.
The isomorphism in (K33I.6) is the canonical cross-product map; the
complex comparison is used only to verify it, so the map is
Galois-equivariant and compatible with trace maps.

For \(G=G_{\mathbf Q(\mu_m),S}\), the odd-prime cohomological-dimension
theorem gives \(\operatorname{cd}_pG=2\), since \(S\) contains the primes
above \(p\) and the archimedean places.  In the Hochschild--Serre
spectral sequence \(H^a(G,H^b_{\rm geom})\), total degree \(3\) has
filtration \(F^1=H^3\), because \(H^3_{\rm geom}=0\).
There is no incoming differential at \((1,2)\), and every outgoing
one has first index at least \(3\).  Consequently there is a
canonical quotient edge
\[
 H^3_{\rm et}(\text{arithmetic product},{\cal F})
   \longrightarrow
 H^1\bigl(G,H^2_{\rm et}(\text{geometric product},{\cal F})\bigr).
                                                               \tag{K33I.7}
\]
The possible term \(H^2(G,H^1_{\rm geom})\) belongs to its kernel;
it introduces no choice of a splitting.

All finite-level groups with finite coefficients are finite.
Consequently the countable inverse systems are Mittag--Leffler,
and the same edge maps pass to the continuous coefficient tower.
Taking the cofinal diagonal in the two internal levels and in
coefficient precision gives
\[
 z_m^{[j]}\longmapsto
 \widetilde z_m^{[j]}
 \in H^1\bigl(G,\,(I_j\widehat\otimes_{\mathcal O}I_j)(-j)\bigr).
                                                               \tag{K33I.8}
\]
One may justify the completed tensor in (K33I.8) by replacing each
finite-coefficient inverse system by its eventually stable images.
This replacement is a pro-isomorphism: at each fixed level, a
sufficiently late transition factors through its stable image.
Tensor these pro-isomorphisms, rather than assuming that an inclusion
of stable images remains injective after tensoring.  The resulting
surjective systems compute the same limit, and their finite tensor
products compute the completed tensor.  Alternatively,
the maps to the finite-dimensional target below are obtained directly
from the finite-level compactly supported dual classes in AC-K33;
every reduction of each chosen dual vector has a finite-level
representative.  Both constructions agree by (K33I.6).
Exactness of inverse limits of the compact modules involved also
follows directly from the finite-intersection property for their
nonempty compact lifting fibres.  These arguments do not assume
surjectivity of every original trace map.

The first coefficient map is the literal KLZ moment and divided-power
multiplication
\[
 {\cal L}^{[j]}(1)
 \xrightarrow{\ \mathrm{mom}^{\,2e-1-j}\cdot\mathrm{id}\ }
 \operatorname{TSym}^{2e-1}{\cal H}(1),            \tag{K33I.9}
\]
followed on cohomology by the fixed eigenform quotient.
It is defined since \(2e-1-j=e\geq0\).
Definition~2.8.1 of KLZ defines this quotient as \(M(f)^*\), with
\(T_\ell'\) and \(U_\ell'\) acting by \(a_\ell(f)\).
In the arithmetic-Frobenius convention its characteristic polynomial
is exactly (K33I.1), so
\[
 M(f)^*=\rho_f,\quad\text{with this declared convention}.         \tag{K33I.10}
\]
For the local normalization, KLZ Theorem~7.2.3(iii--iv), printed
pp.~35--36, gives arithmetic Frobenius \(U_p'\) on the unramified
ordinary quotient and the power \(\kappa^{1+k}\) on the ordinary
submodule; Remark~8.2.9 gives the same normalization after
specialization.  Here \(k=2e-1\).  The quadratic nebentypus
\(\eta_K=\eta_K^{-1}\) eliminates any reciprocal-nebentypus ambiguity.
We use neither the opposite \(M(f)\) convention nor an unspecified
Tate twist.

The second map is the simultaneous quotient already proved in
AC-K33:
\[
 I_j\otimes E \longrightarrow \rho_g^\vee\otimes_E D^\vee.
                                                               \tag{K33I.11}
\]
Its inverse-diamond character is \(u^{-j-1}\), as proved from the
contragredient pairing there, and \(U_p'=0\).
Applying (K33I.9)--(K33I.11) to (K33I.8) gives (K33I.2), with exactly
the remaining twist \(-j\).  Evaluation against \(d\in D\) gives
\(c_{m,d}\).

For completeness, the identification with the actual target is
\[
 \rho_f\otimes\rho_g^\vee(1-e)
       =V\otimes\rho_g^\vee(1)
       \simeq(V\otimes\rho_g)^\vee(1).             \tag{K33I.12}
\]
Indeed the determinant pairing gives \(V^\vee\simeq V\otimes\eta_K\).
For an induced representation, choose the standard two-coset model:
the diagonal map with entries \(1,-1\) intertwines
\(\rho_g^\vee\otimes\eta_K\) with \(\rho_g^\vee\).
These explicitly chosen isomorphisms prove (K33I.12).
In the cubic case \(\chi^c=\chi^{-1}\), so \(\rho_g\) is also
self-dual; this extra fact is compatible with, but is not needed
in place of, the displayed intertwining calculation.
For simultaneous quadratic compensation both \(\tau\) factors
cancel, and the same calculation applies.

We spell out this last coefficient's relation to the original \(K\)-sector.
Write
\[
 U_K=(V|_{G_K}\otimes\chi^b)^\vee(1).
\]
The tensor-induction identity and duality for finite induction give
\[
 U_{\mathbf Q}=(V\otimes\operatorname{Ind}_{G_K}^{G_{\mathbf Q}}
                         \chi^b)^\vee(1)
       \simeq\operatorname{Ind}_{G_K}^{G_{\mathbf Q}}U_K.
                                                               \tag{K33I.13}
\]
These identities are identities of actual representations, not only
of their Euler factors.  We chose \(3\mid N\); since \((m,N)=1\),
the field \(\mathbf Q(\mu_m)\) is unramified at \(3\), whereas \(K\)
is ramified there.  Thus \(K\cap\mathbf Q(\mu_m)=\mathbf Q\).
Restriction of (K33I.13) and Shapiro now give
\[
 H^1(\mathbf Q(\mu_m),U_{\mathbf Q})
       \simeq H^1(K(\mu_m),U_K).                 \tag{K33I.14}
\]
At the base \(m=1\), the prime \(269\) is inert in \(K\), so there is
one local field \(K_p/\mathbf Q_p\), unramified of degree \(2\), and
\[
 H^1(\mathbf Q_p,U_{\mathbf Q})
       \simeq H^1(K_p,U_K).                      \tag{K33I.15}
\]
The restriction map from global to local cohomology commutes with
these Shapiro identifications.  Tensoring finite induction with
\(B_{\rm cris}\) and applying the same Shapiro identity shows that
the kernels defining \(H_f^1\) correspond.  Hence their singular
quotients also correspond; their \(E\)-dimension is the already
computed \(2\) in AC-G69-local.  For general \(m\), local Shapiro
uses the sum over places of \(K(\mu_m)\) above a chosen place of
\(\mathbf Q(\mu_m)\).  We do not assert that the unramified quadratic
extension remains a field after every such local base change.
In particular, (K33I.14) yields the cyclotomic composita
\(K(\mu_m)\), not all \(K\)-ray class fields.

Finally, each coefficient map is bounded on the fixed compact
lattice.  The fixed algebraic moment is integral; the fixed
eigenform projection and AC-K33 map require at most one fixed
power of \(p\) to reach chosen target lattices.  All transport and
corestriction maps of AC-K33-cyclotomic are integral before these
fixed projections.  The same fixed denominator therefore works
for every \(m\).  Functoriality of Galois cohomology makes all the
resulting corestriction identities commute with these maps.
None of these operations proves that the resulting class, or
its local singular determinant, is nonzero.
\end{proof}

\begin{remark}[Source and scope]
The finite-level Hochschild--Serre construction is already described
in KLZ, Definition~3.3.1 and its following paragraph, and the
preordinary tensor edge is explicitly described in its introduction,
printed pp.~5--6.  The proof above supplies the finite-coefficient
and compact-limit details needed for the different, nonordinary
quotient of AC-K33.  Etale comparison, Kunneth, arithmetic
cohomological dimension and finite-coefficient finiteness are
retained standard inputs, not claims of new upstream proofs.
The two de Rham labels in (K33I.11) do not supply two independent
regulator values: that requires an actual Rankin reciprocity
identity on these coefficient maps.
\end{remark}
